\chapter{Conexidad y continuidad}\label{chap:17}

\section{La imposibilidad de separar un espacio}
La compacidad controla cubiertas; la conexidad controla separaciones. Un espacio conexo no puede dividirse en dos regiones abiertas disjuntas y no vacías. Esta propiedad explica el teorema del valor intermedio y se conserva bajo aplicaciones continuas, aunque no implique todavía que dos puntos puedan unirse mediante una curva. Usaremos topologías relativas, cerraduras y continuidad. Este capítulo desarrolla la lectura 17 de \cite{valencia}.

\section{Definiciones equivalentes y ejemplos}
\begin{proposicion}\label{ch17:def}
Para un espacio métrico $X$ son equivalentes:
\begin{enumerate}[label=(\roman*)]
\item los únicos subconjuntos simultáneamente abiertos y cerrados son $\varnothing$ y $X$;
\item no existen abiertos disjuntos no vacíos $U,V$ con $X=U\cup V$;
\item si $U,V$ son abiertos disjuntos y cubren $X$, uno de ellos es vacío.
\end{enumerate}
Si se cumplen, $X$ es \emph{conexo}; en caso contrario es \emph{disconexo}. Una pareja $U,V$ como en (ii) se llama una \emph{separación}. Un subconjunto es conexo si lo es con su topología relativa.\index{conexidad}\index{separacion@separación}
\end{proposicion}
\begin{proof}
Un subconjunto $A$ abierto y cerrado, no vacío y distinto de $X$, produce la separación $A,X\setminus A$: ambos son abiertos, disjuntos y no vacíos. En sentido inverso, en una separación $U,V$ tenemos $V=X\setminus U$, de modo que $U$ es abierto y cerrado. Las condiciones (ii) y (iii) son la misma afirmación expresada por negación y por implicación, respectivamente.
\end{proof}
El vacío y los singletons son conexos. Un espacio discreto no vacío es conexo si y solo si consta de un punto: si tiene al menos dos, cualquier singleton es abierto y cerrado propio.

\begin{lema}\label{ch17:intervalos}
Todo intervalo de $\R$ es conexo; se admiten intervalos abiertos, cerrados, semiabiertos, ilimitados, singletons y el vacío.
\end{lema}
\begin{proof}
Sea $I$ un intervalo y supongamos que tiene una separación $U,V$. Elegimos $a\in U$, $b\in V$; intercambiando los nombres si hace falta, podemos suponer $a<b$. El conjunto $E=U\cap[a,b]$ es no vacío y acotado; sea $c=\sup E$. Por ser $I$ intervalo, $[a,b]\subset I$ y $c\in I$. Existe una sucesión de puntos de $E$ que converge a $c$, por la propiedad del supremo; como $U$ es cerrado en $I$, tenemos $c\in U$. En particular $c\ne b$. Como $U$ es abierto relativo, contiene $I\cap(c-r,c+r)$ para algún $r>0$. Elegimos $t$ con $c<t<\min\{c+r,b\}$. Entonces $t\in E$, contradiciendo la definición de $c$. No existe tal separación.
\end{proof}

\begin{proposicion}\label{ch17:ambiental}
Un subconjunto $A\subset X$ es disconexo si y solo si existen abiertos disjuntos $U,V$ del ambiente $X$ tales que $A\subset U\cup V$ y ambos $A\cap U$, $A\cap V$ son no vacíos.
\end{proposicion}
\begin{proof}
La suficiencia se obtiene intersectando con $A$: aparecen dos abiertos relativos que lo separan. Para la necesidad, sea $A=A_1\cup A_2$ una separación relativa. Cada punto de $A_1$ tiene una bola que no toca $A_2$, porque $A_1$ es abierto en $A$; por ello $\dist(a,A_2)>0$ para $a\in A_1$. La afirmación correspondiente vale intercambiando los índices. Definimos
\[
U=\{x:\dist(x,A_1)<\dist(x,A_2)\},\qquad
V=\{x:\dist(x,A_2)<\dist(x,A_1)\}.
\]
Las distancias a conjuntos no vacíos son funciones continuas, así que $U,V$ son abiertos y disjuntos. En $A_1$ la primera distancia es cero y la segunda positiva; luego $A_1\subset U$. Del mismo modo $A_2\subset V$. Así los abiertos ambientales tienen todas las propiedades requeridas. Es útil construirlos explícitamente: los abiertos ambientales que representan inicialmente una separación relativa no tienen por qué ser disjuntos fuera de $A$.
\end{proof}

\section{Imágenes y familias de conjuntos conexos}
\begin{teorema}\label{ch17:imagen}
La imagen continua de un espacio conexo es conexa.
\end{teorema}
\begin{proof}
Si $f(X)$ tuviera una separación relativa $U,V$, las preimágenes $f^{-1}(U)$ y $f^{-1}(V)$ serían abiertas, disjuntas y cubrirían $X$. Serían no vacías porque ambos miembros de la separación contienen puntos de la imagen. Esto contradice la conexidad de $X$.
\end{proof}

\begin{corolario}\label{ch17:binario}
$X$ es conexo si y solo si toda aplicación continua $X\to\{0,1\}$ es constante. Equivalentemente, toda aplicación continua de $X$ a cualquier espacio discreto es constante.
\end{corolario}
\begin{proof}
Si $X$ es conexo, su imagen continua es conexa y una imagen discreta no puede tener dos puntos. Para la otra dirección, una separación $U,V$ define la función que vale $0$ en $U$ y $1$ en $V$. Es continua porque las preimágenes de los singletons son abiertas, y no es constante. El mismo argumento sirve con dos puntos distintos de cualquier espacio discreto de al menos dos elementos. Para el dominio vacío, usamos la convención habitual de que la función vacía es constante en el sentido de que no toma dos valores diferentes.
\end{proof}

\begin{definicion}
Una familia no vacía $\mathcal C$ de conjuntos no vacíos es \emph{concatenada}\index{familia concatenada} si para cualesquiera $A,B\in\mathcal C$ existe una cadena finita $A=C_1,\ldots,C_m=B$ de miembros de la familia con $C_i\cap C_{i+1}\ne\varnothing$.
\end{definicion}

\begin{proposicion}\label{ch17:union}
La unión de una familia concatenada de conjuntos conexos es conexa. En particular, una familia de conjuntos conexos con un punto común tiene unión conexa.
\end{proposicion}
\begin{proof}
Sea $f:\bigcup\mathcal C\to\{0,1\}$ continua. En cada miembro de $\mathcal C$ es constante, por el corolario \ref{ch17:binario}. Si dos miembros se intersectan, sus valores constantes coinciden al evaluar $f$ en un punto de intersección. Una cadena finita transmite la igualdad entre cualesquiera dos miembros. Así $f$ es constante en toda la unión, que es conexa por el mismo corolario. Si hay un punto común, cualquier pareja de miembros se intersecta y la familia es concatenada. Si una familia contiene conjuntos vacíos, se pueden eliminar sin alterar la unión.
\end{proof}

\begin{teorema}\label{ch17:producto}
El producto finito de espacios métricos no vacíos $X_1,\ldots,X_n$, con su topología producto, es conexo si y solo si cada factor es conexo.
\end{teorema}
\begin{proof}
Si el producto es conexo, la proyección sobre cada factor es continua y sobreyectiva; la imagen es conexa. La sobreyectividad requiere que los demás factores sean no vacíos.

Para la otra dirección, supongamos que todos los factores son conexos. Para $a=(a_1,\ldots,a_n)$ y $i$ definimos la fibra
\[
L_{a,i}=\{x\in\textstyle\prod_jX_j:x_j=a_j\text{ para todo }j\ne i\}.
\]
Aquí se exige igualdad en las coordenadas que no varían; escribir desigualdad cambiaría el conjunto y no permitiría identificarlo con $X_i$. La evaluación de la coordenada $i$ es un homeomorfismo de $L_{a,i}$ sobre $X_i$, con inversa que inserta esa coordenada y deja fijas las demás. Por tanto cada fibra es conexa.

Las fibras cubren el producto y forman una familia concatenada. En efecto, dados puntos $a,b$, las etapas
\[
(a_1,\ldots,a_n),\ (b_1,a_2,\ldots,a_n),\ldots,(b_1,\ldots,b_n)
\]
cambian una coordenada por vez. Cada etapa consecutiva está en una misma fibra, y las fibras sucesivas comparten el punto de transición. Para conectar dos fibras arbitrarias, tomamos un punto en cada una y añadimos al comienzo y al final esas dos fibras a la cadena construida. Aplicamos la proposición \ref{ch17:union}.
\end{proof}

\begin{ejemplo}[La intersección puede romper la conexidad]
Las semicircunferencias superior e inferior de la circunferencia unitaria son conexas: son imágenes continuas de $[0,\pi]$ y $[\pi,2\pi]$ por $t\mapsto(\cos t,\sin t)$. Su intersección es $\{(-1,0),(1,0)\}$, que es disconexa. Por tanto no hay una regla general de estabilidad de la conexidad bajo intersecciones.
\end{ejemplo}

\section{Ejercicios resueltos}
\begin{ejercicio}[Lectura 17, ejercicio 1: precisión del enunciado]
Probar que $X$ es conexo si y solo si todo subconjunto no vacío y \emph{propio} $A\subsetneq X$ tiene frontera no vacía.
\end{ejercicio}
\begin{solucion}
La palabra propio es necesaria, pues $\partial X=\varnothing$ en cualquier espacio. Recordemos $\partial A=\cl A\setminus\interior A$. Si esta frontera es vacía, de $\interior A\subset A\subset\cl A$ se deduce $\interior A=A=\cl A$; por tanto $A$ es abierto y cerrado. Recíprocamente, si $A$ es abierto y cerrado, esas tres expresiones coinciden y su frontera es vacía. Así la condición del ejercicio excluye exactamente los abiertos y cerrados no triviales, que caracterizan la desconexión.
\end{solucion}

\begin{ejercicio}[Lectura 17, ejercicio 2]
Probar: (i) los subconjuntos conexos de $\R$ son exactamente los intervalos; (ii) una función real continua sobre un espacio conexo toma todos los valores intermedios entre dos valores de su imagen; (iii) si toma un valor negativo y uno positivo, tiene un cero.
\end{ejercicio}
\begin{solucion}
Para (i), los intervalos son conexos por el lema \ref{ch17:intervalos}. Si $A\subset\R$ no es un intervalo, hay $a,b\in A$ y $c\notin A$ con $a<c<b$. Entonces $A\cap(-\infty,c)$ y $A\cap(c,\infty)$ son dos abiertos relativos disjuntos no vacíos que cubren $A$, de modo que $A$ no es conexo.

Para (ii), si $X$ es conexo y $f:X\to\R$ continua, el teorema \ref{ch17:imagen} muestra que $f(X)$ es conexo, y por (i) es un intervalo. Si $f(x)=a\le b=f(y)$, todo $t\in[a,b]$ pertenece a $f(X)$, lo que significa que existe $z\in X$ con $f(z)=t$. Para (iii), aplicamos (ii) a $t=0$, que está entre el valor negativo y el positivo. Los extremos de $X$ no necesitan existir: el argumento solo exige dos puntos donde se conozcan los signos.\index{teorema del valor intermedio}\index{teorema de Bolzano}
\end{solucion}

\begin{ejercicio}[Lectura 17, ejercicio 3]
Si $A$ es conexo y $A\subset B\subset\cl A$, probar que $B$ es conexo. Deducir que la cerradura de un conexo es conexa.
\end{ejercicio}
\begin{solucion}
Sea $f:B\to\{0,1\}$ continua. Su restricción a $A$ es constante, digamos de valor $0$. Si $b\in B$ tuviera $f(b)=1$, la preimagen $f^{-1}(\{1\})$ sería una vecindad abierta relativa de $b$ que no toca $A$. Pero $b\in\cl A$ y $A\subset B$, de modo que toda vecindad relativa de $b$ en $B$ toca $A$, contradicción. Así $f$ es constante y $B$ es conexo. Para $A=\varnothing$, también $B=\varnothing$ y el resultado es válido. Tomar $B=\cl A$ da la conclusión final.
\end{solucion}

\begin{ejercicio}[Lectura 17, ejercicio 4]
Probar que todo espacio normado es conexo.
\end{ejercicio}
\begin{solucion}
Para $x\in V$, el segmento $[0,x]=\{tx:0\le t\le1\}$ es imagen continua de $[0,1]$: $\norm{tx-sx}=|t-s|\norm{x}$. Es por tanto conexo. Todos estos segmentos contienen $0$, y su unión es $V$. La proposición \ref{ch17:union} prueba que $V$ es conexo. No se requiere completitud; el argumento vale también en un espacio normado complejo usando el parámetro real $t$.
\end{solucion}

\begin{ejercicio}[Lectura 17, ejercicio 5]
Sea $X$ conexo, $X=A\cup B$ y $A\cap B$ conexo. Probar que $A,B$ son conexos si (i) ambos son cerrados o (ii) ambos son abiertos. En (iii), decidir qué ocurre si $A$ es abierto y $B$ cerrado.
\end{ejercicio}
\begin{solucion}
Si $A\cap B=\varnothing$ y ambos son abiertos, forman una separación salvo que uno sea vacío. Si ambos son cerrados, cada uno es complemento del otro y por tanto también abierto; la misma conclusión se aplica. El miembro no vacío es entonces $X$ y ambos son conexos.

Supongamos ahora $A\cap B\ne\varnothing$. Si $A$ fuera disconexo, existiría una función continua no constante $g:A\to\{0,1\}$. Su restricción al conexo $A\cap B$ sería constante, de valor $c$. Definimos $f=g$ en $A$ y $f=c$ en $B$; las definiciones coinciden en la intersección. Si $A,B$ son abiertos, las preimágenes de abiertos por $f$ son uniones de conjuntos abiertos en $X$, así que $f$ es continua. Si $A,B$ son cerrados, para cada cerrado $F$ del codominio la preimagen es la unión de un cerrado relativo de $A$ y uno de $B$; ambos son cerrados en $X$, luego la unión lo es y $f$ también es continua. En cualquiera de los casos $f$ sería continua no constante sobre $X$, contradicción. Así $A$ es conexo. Intercambiando los papeles de $A$ y $B$, el mismo argumento de pegado prueba que $B$ es conexo.

Para (iii) no se garantiza la conexidad de ninguno de los dos por separado. En $X=\R$, tomemos $A=\R\setminus\{0\}$ y $B=[-1,0]$. El primero es abierto y disconexo, el segundo cerrado, su unión es $\R$ y su intersección $[-1,0)$ es conexa. Para que falle la conexidad de $B$, tomemos $X=[0,2]$, $A=[0,2)$, abierto relativo, y $B=\{1,2\}$, cerrado relativo. Cubren $X$ y su intersección es $\{1\}$, pero $B$ es disconexo. La condición uniforme abierto/abierto o cerrado/cerrado es la que permite el pegado anterior.
\end{solucion}

\begin{ejercicio}[Lectura 17, ejercicio 6]\label{ch17:ex6}
Sean $X$ conexo, $A\subset X$ conexo y $B\subset X\setminus A$ abierto y cerrado en $X\setminus A$. Probar que $A\cup B$ es conexo.
\end{ejercicio}
\begin{solucion}
Si $A$ es vacío, $B$ es abierto y cerrado en $X$ y por tanto es vacío o todo $X$, ambos conexos. Si $A\ne\varnothing$, sea $f:A\cup B\to\{0,1\}$ continua. En $A$ es constante; cambiando los nombres de $0$ y $1$ podemos suponer que vale $0$. Pongamos $D=f^{-1}(\{1\})\subset B$. Este conjunto es abierto y cerrado en $A\cup B$.

Demostraremos que $D$ es abierto y cerrado en $X$. Como $B$ es cerrado en $X\setminus A$, tenemos $\cl B\subset A\cup B$. Por ello $\cl D\subset A\cup B$ y, puesto que $D$ es cerrado en este subespacio, $\cl D=D$: es cerrado en $X$. Para la apertura, si $x\in D$, existe un abierto $O$ de $X$ con $x\in O$ y $O\cap(A\cup B)\subset D$. Como $B$ es abierto en $X\setminus A$, existe un abierto $P$ de $X$ con $x\in P$ y $P\cap(X\setminus A)\subset B$. Todo punto de $O\cap P$ que estuviera en $A$ o en $B$ estaría en $D$ por la primera condición; todo punto que no estuviera en $A$ estaría en $B$ por la segunda. Así $O\cap P\subset D$, y $D$ es abierto. La conexidad de $X$ fuerza $D=\varnothing$ o $D=X$; la segunda posibilidad es imposible porque $A\ne\varnothing$ y $D\cap A=\varnothing$. Entonces $f$ es constante y $A\cup B$ es conexo.

La sugerencia de la lectura es una versión de pegado que también podemos justificar. Si $f$ es continua por separado en $E$ y $F$, y $E\setminus F$ es abierto y cerrado en la diferencia simétrica $E\mathbin{\triangle}F$, entonces $f$ es continua en $E\cup F$, siempre que las definiciones coincidan en la intersección. Un punto de $E\setminus F$ tiene una vecindad relativa en la unión contenida en $E$, por la apertura en la diferencia simétrica; la cerradura da la afirmación correspondiente para $F\setminus E$. En un punto de $E\cap F$, intersectamos las dos vecindades proporcionadas por la continuidad de las restricciones. En los tres casos obtenemos la continuidad en la unión.
\end{solucion}

\begin{ejercicio}[Lectura 17, ejercicio 7]
Construir: (i) un espacio métrico conexo con una bola disconexa; (ii) un conjunto conexo cuyo interior sea disconexo.
\end{ejercicio}
\begin{solucion}
Para (i), en $\R^2$ con la distancia euclídea consideremos
\[
X=(\{-1,1\}\times[0,2])\cup([-1,1]\times\{2\}).
\]
Es la unión de tres segmentos conexos concatenados, luego es conexo. La bola relativa de centro $(-1,0)$ y radio $\sqrt5$ contiene el brazo izquierdo completo, el tramo superior con $-1\le x<0$, y el tramo derecho $\{1\}\times[0,1)$. No contiene ningún punto con primera coordenada entre $0$ y $1$. Sus intersecciones con los semiplanos $x<1/2$ y $x>1/2$ son abiertos relativos disjuntos no vacíos que la cubren. Por tanto la bola es disconexa.

Para (ii), tomemos la unión $K$ de los discos cerrados de radio $1$ centrados en $(-1,0)$ y $(1,0)$. Cada disco es conexo por convexidad, y ambos se tocan en $(0,0)$; su unión es conexa. Su interior es exactamente la unión de los dos discos abiertos, que son disjuntos. Los puntos de frontera distintos del punto de tangencia no son interiores porque tienen vecindades pequeñas que no tocan el otro disco y contienen puntos exteriores al primero. El punto de tangencia tampoco es interior, pues $(0,t)\notin K$ para todo $t\ne0$. Así el interior tiene dos partes abiertas no vacías separadas por el eje vertical y es disconexo.
\end{solucion}

\section{Idea que conviene conservar}
La prueba de conexidad puede reducirse a excluir funciones continuas no constantes con valores $0$ y $1$. Esta perspectiva hace transparentes las imágenes continuas, las uniones concatenadas y los argumentos de pegado. La conexidad pertenece a la topología del espacio, mientras que la forma de sus bolas depende de la métrica elegida.
